%=============================================================================
% TIME CAPSULE – DIGITAL LEGACY APP
% Final Year B.Tech Project Research Paper
% LaTeX Source — IEEE-style Technical Report
%=============================================================================

\documentclass[12pt, a4paper]{article}

%-----------------------------------------------------------------------------
% PACKAGES
%-----------------------------------------------------------------------------
\usepackage[a4paper, top=2.5cm, bottom=2.5cm, left=3cm, right=2.5cm]{geometry}
\usepackage{times}                    % Times New Roman font (professional look)
\usepackage[T1]{fontenc}
\usepackage[utf8]{inputenc}
\usepackage{microtype}                % Better text justification
\usepackage{setspace}
\usepackage{parskip}                  % Space between paragraphs

% Headings and structure
\usepackage{titlesec}
\usepackage{titletoc}
\usepackage{fancyhdr}

% Math
\usepackage{amsmath}
\usepackage{amssymb}

% Tables and figures
\usepackage{booktabs}
\usepackage{longtable}
\usepackage{array}
\usepackage{tabularx}
\usepackage{multirow}
\usepackage{graphicx}
\usepackage{float}
\usepackage{caption}
\usepackage{subcaption}

% Hyperlinks and references
\usepackage[
  colorlinks=true,
  linkcolor=black,
  citecolor=black,
  urlcolor=blue,
  breaklinks=true
]{hyperref}
\usepackage{url}
\usepackage{cite}

% Acronyms and glossary
\usepackage{acronym}

% Appendix
\usepackage[toc,page]{appendix}

% List formatting
\usepackage{enumitem}

% Line spacing
\onehalfspacing

%-----------------------------------------------------------------------------
% HEADER / FOOTER
%-----------------------------------------------------------------------------
\pagestyle{fancy}
\fancyhf{}
\fancyhead[L]{\small Time Capsule -- Digital Legacy Application}
\fancyhead[R]{\small Final Year Project Report}
\fancyfoot[C]{\thepage}
\renewcommand{\headrulewidth}{0.4pt}

%-----------------------------------------------------------------------------
% TITLE INFORMATION
%-----------------------------------------------------------------------------
\title{
  \vspace{-1cm}
  {\LARGE \textbf{Time Capsule: A Secure, Web-Based Digital Legacy Application\\
  with Heartbeat-Triggered Release, Guardian Verification, and Rich Media Composition}}\\[0.5cm]
  {\large Final Year B.Tech Project Report}\\[0.3cm]
  {\normalsize Department of Computer Science and Engineering}\\
  {\normalsize [ADD INSTITUTION NAME]}
}

\author{
  \textbf{[ADD YOUR FULL NAME]}\\
  Roll No.: [ADD ROLL NUMBER]\\
  \texttt{[ADD EMAIL]}\\[0.3cm]
  \textit{Supervised by:}\\
  \textbf{[ADD SUPERVISOR NAME]}\\
  [ADD SUPERVISOR DESIGNATION]
}

\date{April 2026}

%=============================================================================
% DOCUMENT BEGINS
%=============================================================================
\begin{document}

\maketitle
\thispagestyle{empty}

\newpage
\pagenumbering{roman}

%-----------------------------------------------------------------------------
% DECLARATION
%-----------------------------------------------------------------------------
\section*{Declaration}
\addcontentsline{toc}{section}{Declaration}

I hereby declare that the project work titled \textit{``Time Capsule: A Secure, Web-Based Digital Legacy Application with Heartbeat-Triggered Capsule Release''} submitted to [ADD INSTITUTION NAME] in partial fulfilment of the requirements for the award of the degree of Bachelor of Technology in Computer Science and Engineering is a record of original work carried out by me under the guidance of [ADD SUPERVISOR NAME], [ADD DESIGNATION], Department of Computer Science and Engineering.

The results embodied in this report have not been submitted to any other university or institution for the award of any degree or diploma.

\vspace{2cm}

\noindent
\begin{tabular}{p{7cm} p{7cm}}
  \textbf{[ADD YOUR FULL NAME]} & \textbf{[ADD SUPERVISOR NAME]}\\
  Roll No.: [ADD ROLL NUMBER] & [ADD DESIGNATION]\\
  Date: \today & Department of CSE\\
\end{tabular}

\newpage

%-----------------------------------------------------------------------------
% ACKNOWLEDGEMENTS
%-----------------------------------------------------------------------------
\section*{Acknowledgements}
\addcontentsline{toc}{section}{Acknowledgements}

The author wishes to express sincere gratitude to [ADD SUPERVISOR NAME] for the continuous guidance, technical insight, and academic support provided throughout the duration of this project. Appreciation is also extended to the faculty members of the Department of Computer Science and Engineering at [ADD INSTITUTION NAME] for their encouragement and constructive feedback during review sessions.

The author acknowledges the open-source communities behind the Flask, React, SQLAlchemy, and Python Cryptography libraries, whose well-documented, freely available tools made the implementation of this system feasible within an academic timeline.

\newpage

%-----------------------------------------------------------------------------
% ABSTRACT
%-----------------------------------------------------------------------------
\begin{abstract}
\addcontentsline{toc}{section}{Abstract}

\noindent The exponential growth of personal digital data has created a critical yet largely unaddressed problem: the secure and intentional management of digital assets after death or prolonged incapacitation. Existing solutions such as password managers and generic cloud storage lack dedicated mechanisms for privacy-preserving, time-deferred delivery of personal communications and digital legacy materials. This paper presents \textit{Time Capsule}, a full-stack, web-based digital legacy application designed to enable users to create, encrypt, and schedule the future delivery of personal messages, files, and other digital artefacts to designated recipients.

The system employs a decoupled architecture comprising a React-based single-page application (SPA) frontend and a Python Flask REST API backend, with PostgreSQL as the persistence layer managed through the SQLAlchemy Object Relational Mapper (ORM). A central security contribution of the system is the application of Fernet symmetric encryption — built upon the Advanced Encryption Standard (AES-128-CBC) — to ensure confidentiality of capsule contents at rest. Authentication is achieved through JSON Web Tokens (JWT) with server-side validation, and passwords are stored using Werkzeug's PBKDF2-SHA256 hashing scheme.

The primary novel contribution of this project is the design and implementation of a \textit{heartbeat-triggered capsule release mechanism}: a background monitoring system that periodically queries user liveness through automated email pings. If no confirmation is received within a configurable inactivity window, designated capsules are automatically released to their intended recipients. This mechanism addresses the scenario of unexpected user incapacitation and distinguishes this system from purely time-based scheduled delivery platforms.

Beyond the core platform, two extended features were designed and implemented as part of this project. First, a \textit{Guardian Verification Portal}: a tokenised, public-facing web interface through which designated trusted contacts can formally confirm or deny event-based capsule release requests, with a complete audit log maintained per capsule. Capsule delivery is triggered automatically when all guardians confirm. Second, \textit{Rich Text and Media Composition}: a WYSIWYG message editor built on the TipTap framework, enabling formatted text input, and an in-browser video recording module using the MediaRecorder API, allowing users to attach recorded video messages directly to capsules.

Automated scheduling of capsule delivery is achieved through APScheduler, a lightweight Python background job scheduler integrated into the Flask application lifecycle. An SMTP-based email delivery service handles all outbound communications. The paper details the system's requirements, architectural design, security model, implementation methodology, and testing strategy, and evaluates the system's fitness as a practical, privacy-preserving digital legacy platform.

\vspace{0.5cm}
\noindent\textbf{Keywords:} digital legacy, time capsule, heartbeat inactivity detection, Fernet encryption, AES-128, JWT authentication, Flask REST API, React, APScheduler, guardian verification portal, TipTap WYSIWYG, MediaRecorder API, video message, scheduled email delivery, secure web application, posthumous data management.
\end{abstract}

\newpage

%-----------------------------------------------------------------------------
% TABLE OF CONTENTS
%-----------------------------------------------------------------------------
\tableofcontents

\newpage
\listoffigures
\addcontentsline{toc}{section}{List of Figures}

\newpage
\listoftables
\addcontentsline{toc}{section}{List of Tables}

\newpage
\pagenumbering{arabic}

%=============================================================================
% SECTION 1 — INTRODUCTION
%=============================================================================
\section{Introduction}
\label{sec:introduction}

\subsection{Background and Motivation}

The digitisation of personal life has progressed at an unprecedented rate over the past two decades. Individuals today accumulate vast repositories of digital assets: photographs, written correspondence, video recordings, financial records, legal documents, and personal reflections stored across cloud services, social media platforms, and local devices. Yet despite this pervasive digital presence, there exists a profound absence of structured, privacy-preserving mechanisms for the intentional and secure transmission of these assets to loved ones in the event of the asset holder's death or incapacitation.

This condition has been described in archival and information science literature as the onset of a \textit{digital dark age} \cite{dinneen2024} — a paradox where the most information-rich generation in human history risks leaving behind digital estates that are either irretrievably inaccessible or vulnerable to misuse due to the absence of clear inheritance protocols. The problem is compounded by legal ambiguity surrounding digital assets: most jurisdictions lack cohesive statutory frameworks governing the posthumous transfer of digital property \cite{mali2021}, leaving families reliant on platform-specific policies that are inconsistently designed and difficult to navigate under grief.

Beyond the legal dimension, there is a deeply human motivation for the secure, intentional delivery of personal messages across time. Letters addressed to children on milestone birthdays, farewell messages for loved ones, guidance notes for unforeseen circumstances, or time-sensitive disclosures — these are forms of communication that do not belong to any existing digital service category. Scheduling platforms do not guarantee privacy; note-taking applications do not deliver; and traditional cloud storage does not manage recipients or release timing. The gap between intent and execution is wide, and it is precisely this gap that the present system addresses.

\subsection{Problem Statement}

The following core problems motivate the design and implementation of this project:

\begin{enumerate}[label=\arabic*.]
  \item \textbf{Absence of privacy-preserving delayed delivery:} No widely available consumer platform combines encrypted message storage with automated, time-deferred, recipient-directed email delivery.
  \item \textbf{Lack of inactivity-triggered release:} No open-source or academic system adequately addresses the scenario where a user becomes incapacitated without prior scheduling — the digital equivalent of having no will at all.
  \item \textbf{Inadequate data confidentiality:} Existing digital legacy tools typically lack encryption at rest, exposing sensitive personal communications to potential administrative or adversarial access.
  \item \textbf{Platform fragmentation:} Managing digital legacies currently requires relying on multiple disjoint services with no unified interface for recipients, delivery tracking, or guardian-based verification.
\end{enumerate}

\subsection{Project Objectives}

The Time Capsule project is designed to achieve the following objectives:

\begin{itemize}
  \item Design and implement a secure, web-based digital legacy platform enabling authenticated users to create encrypted personal capsules for future delivery.
  \item Implement a \textit{heartbeat inactivity monitoring system} that detects prolonged user absence and automatically triggers event-based capsule release.
  \item Provide a multi-recipient delivery mechanism through which a single capsule may be addressed to multiple designated individuals.
  \item Enforce \textit{encryption at rest} for all capsule content using AES-128-based Fernet symmetric encryption.
  \item Establish a secure, stateless REST API with JWT-based authentication and CORS-controlled access.
  \item Build an accessible, responsive single-page application frontend to enable capsule management without technical expertise.
  \item Implement attachment support for digital files (photographs, documents, and other media) as part of capsule content.
  \item Design and implement a \textit{Guardian Verification Portal}: a tokenised, public-facing web interface allowing designated trusted contacts to formally confirm or deny event-based capsule release requests, with full audit logging and automatic delivery upon unanimous confirmation.
  \item Integrate a \textit{WYSIWYG rich text editor} (TipTap) for formatted message composition and an \textit{in-browser video recording} module (MediaRecorder API) enabling users to attach recorded video messages to capsules.
\end{itemize}

\subsection{Scope and Limitations}

This project delivers a functionally complete web application encompassing user registration, capsule lifecycle management, encrypted storage, scheduled delivery, and heartbeat monitoring. The system is scoped as a single-user-class platform — that is, all users are creators, and recipients receive capsules through email without requiring their own accounts. The current scope does not include SMS or push notification delivery, client-side end-to-end encryption (messages are encrypted on the server), or blockchain-based audit trails. These are identified as future work. The system targets web-based deployment and is optimised for modern browsers.

\subsection{Report Structure}

The remainder of this report is organised as follows. Section~\ref{sec:literature} reviews related work in digital legacy management, encryption standards, authentication mechanisms, and scheduling systems. Section~\ref{sec:requirements} specifies functional and non-functional system requirements. Section~\ref{sec:design} describes the system architecture, database schema, and process flows. Section~\ref{sec:implementation} details the implementation methodology, technology choices, and key subsystems. Section~\ref{sec:security} analyses the system's security posture and testing strategy. Section~\ref{sec:results} presents results and observations from system operation. Section~\ref{sec:conclusion} concludes the paper and outlines directions for future development.

%=============================================================================
% SECTION 2 — LITERATURE REVIEW
%=============================================================================
\section{Literature Review and Related Work}
\label{sec:literature}

\subsection{Digital Legacy and Posthumous Data Management}

The academic study of digital legacy has grown substantially over the past decade, driven by increased awareness of the social and informational implications of death in digital societies. Massimi and Baecker \cite{massimi2010} were among the earliest researchers to identify the need for dedicated technological support for bereaved families in managing the digital assets of deceased relatives, noting that existing platforms were designed without death in mind. Brubaker et al. \cite{brubaker2014} extended this discourse by examining the social responsibilities that arise from posthumous data stewardship, highlighting the ethical complexity of managing another person's digital identity.

Maciel et al. \cite{maciel2021, maciel2022} proposed a systemic framework for Digital Legacy Management Systems (DLMS), identifying key requirements including content preservation, access control, and designated recipient management — requirements that directly inform the architecture of the present system. Bahri, Carminati, and Ferrari \cite{bahri2015} further explored the problem of posthumous online social estate management, proposing an integrated model that combines legal, technical, and social considerations. Dinneen \cite{dinneen2024}, in a comprehensive literature review, characterised information science's evolving intersection with death, identifying digital legacy as a formative subdiscipline requiring dedicated technical solutions.

Chen et al. \cite{chen2025} proposed a digital will solution for posthumous data management incorporating cryptographic mechanisms, providing technical grounding for the approach taken in this project. The intersection of digital assets with legal frameworks is addressed by Mali and Prakash \cite{mali2021}, who analyse ownership, inheritance, and the legal implications of digital afterlife — a dimension that contextualises the social necessity of the proposed system. Hopkins \cite{hopkins2013} similarly examined the concept of managing a digital estate in the cloud, drawing parallels to traditional estate planning and identifying the need for structured digital succession tools.

\subsection{Related Systems}

\subsubsection{Google Inactive Account Manager}

The most prominent industry implementation of inactivity-triggered data management is Google's \textit{Inactive Account Manager} (IAM), launched in 2013 \cite{google2013blog, google2024docs}. The IAM allows Google account holders to specify trusted contacts who will be notified if the account becomes inactive for a configurable period, ranging from three to eighteen months. Upon triggering, designated contacts may be granted access to specified Google service data, or the entire account may be automatically deleted.

While the IAM represents a significant step toward industry-recognised digital inheritance, it has several limitations in the context of the present work: (i) it is proprietary and tightly coupled to the Google ecosystem; (ii) it does not support encrypted, personalised message composition directed at specific recipients; (iii) it does not permit content scheduling for specific future dates; and (iv) it lacks a capsule-based mental model for curating individual pieces of digital legacy. The Time Capsule system specifically addresses these limitations by providing a platform-agnostic, content-centric, privacy-preserving alternative.

\subsubsection{Scheduled Messaging Platforms}

Several consumer applications — such as FutureMe \cite{futureme}, Dead Man's Switch \cite{deadmansswitch}, and LifeVault — offer variants of scheduled message delivery. FutureMe permits users to email their future selves on a chosen date but lacks multi-recipient delivery, encryption, attachment support, and inactivity triggering. Dead Man's Switch employs an inactivity-based model conceptually similar to the heartbeat mechanism described in this paper, but relies on unencrypted plain-text messages and lacks a structured digital legacy framework. None of these systems implement encryption at rest, JWT-secured APIs, or structured capsule lifecycle management. The present project aims to synthesise and formally engineer these disparate capabilities into a single, cohesive, security-focused platform.

\subsection{Encryption at Rest}

The confidentiality of stored digital data is a well-established security requirement. NIST FIPS 197 \cite{nist2001} specifies the Advanced Encryption Standard (AES) as the approved symmetric encryption algorithm for protecting sensitive government information, and it has since become the de facto global standard for data confidentiality. NIST Special Publication 800-111 \cite{nist2007} provides guidelines specifically for storage encryption technologies, recommending AES-based approaches for end-user data protection.

This project employs the \textit{Fernet} symmetric encryption scheme, implemented via the Python \texttt{cryptography} library \cite{pyca2024}. Fernet is a high-level, authenticated encryption abstraction built atop AES-128-CBC with HMAC-SHA256 for message authentication. The use of authenticated encryption ensures both confidentiality and integrity of stored capsule content — a property essential in a system handling sensitive personal communications. Backendal et al. \cite{backendal2024ieee, backendal2024crypto} provide a formal treatment of end-to-end encrypted cloud storage, establishing the theoretical grounding for the threat model addressed by encryption at rest in this system.

\subsection{Authentication Mechanisms}

JSON Web Tokens (JWT) are specified by IETF RFC 7519 \cite{rfc7519} as a compact, URL-safe means of representing claims between two parties. In web applications, JWTs are widely adopted for stateless authentication: the server signs a token containing the user's identity and expiry time, and the client includes this token in subsequent API requests. This eliminates the need for server-side session storage, improving scalability.

Alshehri et al. \cite{alshehri2023} demonstrated that behaviour-aware JWT implementations can further enhance authentication security by incorporating contextual signals. Fu and Liu \cite{fu2022} provide a critical security analysis of JWT implementations in web applications, identifying common vulnerabilities — including algorithm confusion attacks and improper signature validation — that must be mitigated in production systems. Siriwardena \cite{siriwardena2014} offers a comprehensive treatment of API security using JWT and related standards, informing the authentication design of this project.

\subsection{Web Application Security}

The OWASP API Security Top 10 \cite{owasp2023} identifies the ten most critical API security risks, including broken object-level authorisation, broken authentication, excessive data exposure, and security misconfigurations. The design of the Time Capsule API explicitly addresses several of these risks through owner-level access control on all resources, JWT-based authentication with appropriate expiry, minimal data serialisation in responses, and CORS policy enforcement. Cross-Origin Resource Sharing (CORS) \cite{cors2024} is a W3C mechanism that restricts browser-initiated cross-origin HTTP requests; proper CORS configuration is essential for decoupled SPA-backend architectures and is a security control implemented in this system.

NIST Special Publication 800-63B \cite{nist800-63b} provides authoritative guidelines on digital identity, authentication, and session management, including recommendations for idle session timeout thresholds. These guidelines contextualise the heartbeat inactivity window design described in Section~\ref{sec:design}.

\subsection{Background Scheduling Systems}

Automated, time-based job execution is a foundational requirement of the Time Capsule delivery system. APScheduler \cite{apscheduler} is a mature Python scheduling library supporting in-process scheduling with multiple job store backends, including SQLAlchemy-based persistence. It is well-suited to integration within Flask applications and supports both interval-based and cron-style job triggering. For the present system, an interval-based scheduler checking for due capsules every sixty seconds provides an acceptable trade-off between delivery timeliness and computational overhead.

\subsection{Summary and Positioning}

The literature review establishes that while the problem of digital legacy management has been well-studied from social, ethical, and legal perspectives, technically robust implementations remain sparse — particularly those combining encryption at rest, JWT-secured APIs, multi-recipient delivery, and heartbeat-triggered release in a unified system. The Time Capsule project directly addresses this gap, contributing a practical, secure, and extensible digital legacy platform grounded in established cryptographic and web security standards.

%=============================================================================
% SECTION 3 — SYSTEM REQUIREMENTS AND ANALYSIS
%=============================================================================
\section{System Requirements and Analysis}
\label{sec:requirements}

\subsection{Stakeholder Analysis}

The system involves three primary stakeholder categories:

\begin{itemize}
  \item \textbf{Capsule Creators (Users):} Authenticated individuals who create and manage capsules, define recipients and guardians, and configure delivery rules.
  \item \textbf{Recipients:} Individuals designated to receive capsule deliveries. Recipients do not require platform accounts; they interact with the system solely through delivery emails.
  \item \textbf{Guardians:} Trusted contacts designated by users to verify event-based release triggers. Guardians confirm or deny inactivity-based capsule release requests through email-based confirmation links.
\end{itemize}

\subsection{Functional Requirements}

Table~\ref{tab:functional_req} enumerates the functional requirements of the system, organised by module.

\begin{longtable}{|p{0.8cm}|p{3.5cm}|p{8.5cm}|}
\caption{Functional Requirements}
\label{tab:functional_req}\\
\hline
\textbf{ID} & \textbf{Module} & \textbf{Requirement Description} \\
\hline
\endfirsthead
\multicolumn{3}{c}{\textit{(Continued from previous page)}}\\
\hline
\textbf{ID} & \textbf{Module} & \textbf{Requirement Description} \\
\hline
\endhead
\hline
\multicolumn{3}{r}{\textit{Continued on next page}}\\
\endfoot
\hline
\endlastfoot
FR01 & Authentication & Users shall be able to register with a unique email address and a password meeting minimum complexity requirements. \\
\hline
FR02 & Authentication & Users shall be able to authenticate via email and password and receive a JWT for subsequent API access. \\
\hline
FR03 & Authentication & The system shall reject expired or invalid JWT tokens with a 401 Unauthorised response. \\
\hline
FR04 & Recipients & Authenticated users shall be able to create, view, update, and delete recipient profiles. \\
\hline
FR05 & Recipients & A recipient profile shall store at minimum a name, email address, and optional relationship descriptor. \\
\hline
FR06 & Capsules & Authenticated users shall be able to create a capsule with a title, message body, one or more recipients, and a release configuration. \\
\hline
FR07 & Capsules & Capsule messages shall be encrypted before database persistence and decrypted only upon authenticated retrieval by the owning user. \\
\hline
FR08 & Capsules & Capsules shall support two release types: \texttt{TIME} (specific scheduled datetime) and \texttt{EVENT} (inactivity-triggered). \\
\hline
FR09 & Capsules & The system shall support four capsule lifecycle states: \texttt{DRAFT}, \texttt{SCHEDULED}, \texttt{SENT}, and \texttt{CANCELLED}. \\
\hline
FR10 & Capsules & Users shall be able to edit capsules in \texttt{DRAFT} or \texttt{SCHEDULED} state. \texttt{SENT} and \texttt{CANCELLED} capsules shall be immutable. \\
\hline
FR11 & Attachments & Users shall be able to upload file attachments (images, documents, video) to a capsule, subject to size and MIME-type restrictions. \\
\hline
FR12 & Attachments & The system shall validate uploaded file types against an approved MIME-type allowlist before storage. \\
\hline
FR13 & Delivery & The scheduler shall automatically deliver \texttt{SCHEDULED} capsules whose release datetime has elapsed. \\
\hline
FR14 & Delivery & On delivery, the system shall send a personalised HTML email to each capsule recipient containing the decrypted message and attachment references. \\
\hline
FR15 & Delivery & All delivery attempts shall be recorded in a delivery log with status (\texttt{PENDING}, \texttt{SENT}, \texttt{FAILED}) and timestamp. \\
\hline
FR16 & Heartbeat & The system shall periodically send heartbeat ping emails to all active users. \\
\hline
FR17 & Heartbeat & If a user fails to confirm a heartbeat within the configured inactivity window, the system shall trigger all of that user's event-based capsules. \\
\hline
FR18 & Guardians & Users shall be able to designate trusted guardians for event-based capsules. \\
\hline
FR19 & Dashboard & Authenticated users shall be able to view a dashboard summarising their capsule statuses, recipient count, and recent delivery activity. \\
\hline
FR20 & Guardian Verification & Authenticated users shall be able to request guardian verification for an event-based capsule; the system shall generate unique tokenised links and dispatch them to each assigned guardian via email. \\
\hline
FR21 & Guardian Verification & Guardians shall be able to access a public, token-authenticated web portal to confirm or deny a capsule release request, optionally providing notes; the system shall record the decision, timestamp, and IP address in an audit log. \\
\hline
FR22 & Guardian Verification & When all guardians for a capsule have confirmed, the system shall automatically trigger immediate capsule delivery without waiting for the inactivity window to elapse. \\
\hline
FR23 & Rich Composition & The capsule creation and editing interface shall provide a WYSIWYG rich text editor supporting formatted text (bold, italic, underline, headings, lists, links, blockquotes), with the resulting HTML content encrypted and stored server-side. \\
\hline
FR24 & Rich Composition & The capsule creation and editing interface shall provide an in-browser video recording module allowing users to record, preview, and attach video messages to capsules using the browser's MediaRecorder API. \\
\hline
\end{longtable}

\subsection{Non-Functional Requirements}

\begin{longtable}{|p{0.8cm}|p{2.8cm}|p{9cm}|}
\caption{Non-Functional Requirements}
\label{tab:nonfunctional_req}\\
\hline
\textbf{ID} & \textbf{Category} & \textbf{Requirement Description} \\
\hline
\endfirsthead
\hline
\textbf{ID} & \textbf{Category} & \textbf{Requirement Description} \\
\hline
\endhead
\hline
NFR01 & Security & All capsule message content shall be encrypted at rest using AES-128-CBC (Fernet). \\
\hline
NFR02 & Security & User passwords shall never be stored in plaintext; PBKDF2-SHA256 hashing with Werkzeug shall be applied. \\
\hline
NFR03 & Security & API endpoints shall be protected by JWT-based authentication; unauthenticated access to protected resources shall return HTTP 401. \\
\hline
NFR04 & Security & CORS policy shall restrict API access to whitelisted frontend origins only. \\
\hline
NFR05 & Security & File uploads shall be limited to 16 MB per attachment and restricted to an approved MIME-type allowlist. \\
\hline
NFR06 & Reliability & The background scheduler shall check for deliverable capsules at intervals not exceeding 60 seconds. \\
\hline
NFR07 & Reliability & Email delivery failures shall be logged and shall not cause application crashes; the system shall implement asynchronous email delivery to prevent request blocking. \\
\hline
NFR08 & Performance & API response times for standard CRUD operations shall not exceed 2 seconds under normal load. \\
\hline
NFR09 & Usability & The frontend shall be responsive and usable on desktop and mobile screen sizes without requiring native app installation. \\
\hline
NFR10 & Maintainability & The backend codebase shall use a modular Blueprint-based structure to isolate functional domains. \\
\hline
NFR11 & Portability & The system shall support deployment on any WSGI-compatible Python 3.12+ environment. \\
\hline
NFR12 & Data Integrity & Database schema changes shall be managed through Alembic-based Flask-Migrate versioned migrations. \\
\hline
\end{longtable}

\subsection{Constraints and Assumptions}

The system operates under the following design constraints:
\begin{itemize}
  \item The encryption key is symmetric: the server decrypts messages on behalf of authenticated users. This represents a server-trusted model rather than end-to-end encryption, a deliberate trade-off to support attachment delivery and server-side scheduling.
  \item Email delivery is dependent on SMTP service availability. Transient failures are logged but not automatically retried in the current implementation.
  \item The heartbeat inactivity window is a configurable system parameter; its default value is assumed to be 30 days \textit{[VERIFY with actual configured value]}.
  \item The system does not implement two-factor authentication in its current scope.
\end{itemize}

%=============================================================================
% SECTION 4 — SYSTEM DESIGN AND ARCHITECTURE
%=============================================================================
\section{System Design and Architecture}
\label{sec:design}

\subsection{Architectural Overview}

The Time Capsule system employs a \textit{decoupled three-tier architecture} comprising a presentation tier, an application logic tier, and a data persistence tier. This separation of concerns improves maintainability, enables independent scaling of frontend and backend components, and allows the frontend to be served from a different host than the API — a configuration particularly suited to modern cloud deployment patterns.

\begin{figure}[H]
  \centering
  \fbox{\parbox{0.85\textwidth}{\centering\vspace{2.5cm}[ADD SYSTEM ARCHITECTURE DIAGRAM — showing React SPA, Flask API, PostgreSQL, APScheduler, SMTP service, and their communication channels]\vspace{2.5cm}}}
  \caption{System Architecture Overview}
  \label{fig:architecture}
\end{figure}

The three tiers are:

\begin{enumerate}
  \item \textbf{Presentation Tier:} A React 18 single-page application, bundled by Vite, served statically. The frontend communicates with the backend exclusively through a well-defined REST API over HTTPS, using an Axios HTTP client configured with a JWT Bearer token interceptor.

  \item \textbf{Application Logic Tier:} A Python Flask 3.x REST API that implements all business logic, including user authentication, capsule lifecycle management, encryption and decryption of messages, file handling, and background scheduling. The API is modularised using Flask Blueprints into four functional domains: \texttt{auth}, \texttt{capsules}, \texttt{recipients}, and \texttt{main}.

  \item \textbf{Data Persistence Tier:} A PostgreSQL relational database accessed through the SQLAlchemy ORM. Schema versioning is managed via Flask-Migrate (Alembic). File attachments are stored on the server filesystem under a structured path hierarchy.
\end{enumerate}

Additionally, the system includes two background service components that operate independently of the request-response cycle:

\begin{itemize}
  \item \textbf{APScheduler:} An in-process background scheduler that executes periodic tasks including capsule delivery checking, reminder email dispatch, and heartbeat monitoring.
  \item \textbf{SMTP Email Service:} An asynchronous email dispatch module using Flask-Mail, executed in background threads to avoid blocking API responses.
\end{itemize}

\subsection{Database Design}

The relational data model comprises ten entities designed to represent the full lifecycle of a digital capsule, from creation through delivery.

\begin{figure}[H]
  \centering
  \fbox{\parbox{0.85\textwidth}{\centering\vspace{3.5cm}[ADD ENTITY-RELATIONSHIP DIAGRAM (ERD) — showing all 10 entities: User, Capsule, Recipient, Guardian, CapsuleRecipient, CapsuleGuardian, Attachment, DeliveryLog, HeartbeatCheck, GuardianVerificationRequest with their attributes and relationships]\vspace{3.5cm}}}
  \caption{Entity-Relationship Diagram}
  \label{fig:erd}
\end{figure}

\subsubsection{Core Entities}

\textbf{User:} Represents an authenticated platform member. Key attributes include \texttt{id} (UUID primary key), \texttt{email} (unique, indexed), \texttt{password\_hash} (PBKDF2-SHA256), \texttt{full\_name}, \texttt{created\_at}, and \texttt{last\_heartbeat\_at} (timestamp of the most recent confirmed heartbeat, used by the inactivity monitoring subsystem).

\textbf{Capsule:} The central entity of the data model. Attributes include \texttt{id}, \texttt{owner\_id} (foreign key to User), \texttt{title}, \texttt{encrypted\_message} (ciphertext stored as TEXT), \texttt{release\_type} (\texttt{TIME} or \texttt{EVENT}), \texttt{release\_at} (nullable datetime for time-based release), \texttt{status} (\texttt{DRAFT}, \texttt{SCHEDULED}, \texttt{SENT}, \texttt{CANCELLED}), \texttt{created\_at}, and \texttt{updated\_at}.

\textbf{Recipient:} Stores contact information for individuals who will receive capsules. Attributes include \texttt{id}, \texttt{owner\_id} (foreign key to User), \texttt{name}, \texttt{email}, \texttt{relationship} (e.g., ``daughter'', ``lawyer''), and \texttt{created\_at}.

\textbf{Guardian:} Stores trusted verifier contacts designated by a user. Attributes mirror the Recipient entity with the addition of a \texttt{guardian\_role} descriptor.

\subsubsection{Junction and Auxiliary Entities}

\textbf{CapsuleRecipient:} A many-to-many junction table linking Capsule to Recipient, allowing a single capsule to address multiple recipients.

\textbf{CapsuleGuardian:} A many-to-many junction table linking Capsule to Guardian, enabling guardian-verified event-based releases.

\textbf{Attachment:} Stores metadata for files uploaded to a capsule: \texttt{id}, \texttt{capsule\_id} (foreign key), \texttt{original\_filename}, \texttt{stored\_filename} (UUID-based to prevent collision), \texttt{mime\_type}, \texttt{file\_size}, \texttt{storage\_path}, and \texttt{uploaded\_at}.

\textbf{DeliveryLog:} Records each delivery attempt per recipient: \texttt{id}, \texttt{capsule\_id}, \texttt{recipient\_id}, \texttt{status} (\texttt{PENDING}, \texttt{SENT}, \texttt{FAILED}), \texttt{attempted\_at}, and \texttt{error\_message} (for failure diagnosis).

\textbf{HeartbeatCheck:} Tracks the state of heartbeat pings per user: \texttt{id}, \texttt{user\_id}, \texttt{sent\_at}, \texttt{confirmed\_at} (nullable), \texttt{status} (\texttt{AWAITING}, \texttt{CONFIRMED}, \texttt{EXPIRED}), and \texttt{confirmation\_token} (used in the response link included in heartbeat emails).

\textbf{GuardianVerificationRequest:} Records each formal guardian verification request issued for an event-based capsule. Attributes include \texttt{id}, \texttt{capsule\_id} (foreign key to Capsule, cascade delete), \texttt{guardian\_id} (foreign key to Guardian), \texttt{token} (a 64-character URL-safe random token, unique, indexed), \texttt{status} (\texttt{PENDING}, \texttt{CONFIRMED}, \texttt{DENIED}, \texttt{EXPIRED}), \texttt{sent\_at}, \texttt{responded\_at} (nullable), \texttt{response\_notes} (optional free-text from guardian), and \texttt{ip\_address} (IP of the guardian at response time, for audit purposes).

Table~\ref{tab:db_summary} summarises the ten database models and their roles.

\begin{table}[H]
  \centering
  \caption{Database Entity Summary}
  \label{tab:db_summary}
  \begin{tabularx}{\textwidth}{|l|X|}
  \hline
  \textbf{Entity} & \textbf{Role} \\
  \hline
  User & Platform account holder and capsule creator \\
  \hline
  Capsule & Core unit of digital legacy content \\
  \hline
  Recipient & Designated delivery target (no platform account required) \\
  \hline
  Guardian & Trusted contact for event-based release verification \\
  \hline
  CapsuleRecipient & Many-to-many: Capsule $\leftrightarrow$ Recipient \\
  \hline
  CapsuleGuardian & Many-to-many: Capsule $\leftrightarrow$ Guardian \\
  \hline
  Attachment & File metadata linked to a capsule \\
  \hline
  DeliveryLog & Audit log of all capsule delivery attempts \\
  \hline
  HeartbeatCheck & Inactivity monitoring state per user \\
  \hline
  GuardianVerificationRequest & Tokenised guardian confirm/deny audit record per capsule \\
  \hline
  \end{tabularx}
\end{table}

\subsection{Process Flow Designs}

\subsubsection{Capsule Creation and Encrypted Storage Flow}

Figure~\ref{fig:capsule_creation} illustrates the capsule creation flow. When a user submits a new capsule through the frontend:

\begin{enumerate}
  \item The frontend sends a \texttt{POST /api/capsules} request with the JWT Bearer token in the \texttt{Authorization} header.
  \item The Flask \texttt{@token\_required} decorator validates the JWT signature and expiry; on failure, a 401 response is returned immediately.
  \item The route handler extracts and validates the request payload, verifying that all referenced recipient IDs belong to the authenticated user.
  \item The plaintext message is passed to the \texttt{encrypt()} function of the encryption service, which applies Fernet symmetric encryption using the server-managed key.
  \item The resulting ciphertext is stored in the \texttt{encrypted\_message} field of the Capsule record; the plaintext is never persisted.
  \item \texttt{CapsuleRecipient} junction records are created to link the capsule to its recipients.
  \item A 201 Created response is returned with the capsule metadata (excluding the encrypted message content).
\end{enumerate}

\begin{figure}[H]
  \centering
  \fbox{\parbox{0.85\textwidth}{\centering\vspace{2.5cm}[ADD SEQUENCE DIAGRAM or FLOWCHART — Capsule Creation \& Encrypted Storage Flow]\vspace{2.5cm}}}
  \caption{Capsule Creation and Encrypted Storage Flow}
  \label{fig:capsule_creation}
\end{figure}

\subsubsection{Scheduled Delivery Flow}

Figure~\ref{fig:delivery_flow} illustrates the automated scheduled delivery process. APScheduler executes the delivery job at sixty-second intervals:

\begin{enumerate}
  \item The scheduler queries the database for all capsules where \texttt{status = SCHEDULED} and \texttt{release\_at $\leq$ now()}.
  \item For each qualifying capsule, the encrypted message is retrieved from the database.
  \item The \texttt{decrypt()} function of the encryption service is invoked to recover the plaintext message.
  \item For each linked recipient, an HTML delivery email is composed and dispatched asynchronously via the SMTP service.
  \item A \texttt{DeliveryLog} record is created for each recipient, recording the attempt status.
  \item Upon successful delivery to all recipients, the capsule's \texttt{status} is updated to \texttt{SENT}.
  \item On failure, the delivery log records the error; the capsule remains in \texttt{SCHEDULED} state to be retried in the next scheduler cycle.
\end{enumerate}

\begin{figure}[H]
  \centering
  \fbox{\parbox{0.85\textwidth}{\centering\vspace{2.5cm}[ADD FLOWCHART — Scheduled Capsule Delivery Flow (APScheduler cycle)]\vspace{2.5cm}}}
  \caption{Automated Scheduled Delivery Flow}
  \label{fig:delivery_flow}
\end{figure}

\subsubsection{Heartbeat Inactivity Monitoring Flow}

The heartbeat mechanism is the system's most novel contribution. It addresses the scenario in which a user has not pre-scheduled any capsules but becomes incapacitated. Figure~\ref{fig:heartbeat_flow} illustrates the process:

\begin{enumerate}
  \item The heartbeat scheduler job runs at a configured periodic interval (e.g., every 7 days \textit{[VERIFY]}).
  \item For each active user, the job computes the elapsed time since the user's most recent confirmed heartbeat or account activity.
  \item If the elapsed time exceeds the defined \textit{ping threshold} (e.g., 30 days \textit{[VERIFY]}), a \texttt{HeartbeatCheck} record is created with status \texttt{AWAITING}, and a heartbeat ping email is dispatched to the user.
  \item The email contains a unique, tokenised confirmation link. When clicked, the confirmation endpoint updates the \texttt{HeartbeatCheck} to \texttt{CONFIRMED} and resets the user's \texttt{last\_heartbeat\_at} timestamp.
  \item If the user does not confirm within the \textit{expiry window} (e.g., 30 additional days \textit{[VERIFY]}), the check is marked \texttt{EXPIRED}.
  \item Upon expiry, the scheduler identifies all \texttt{EVENT}-type capsules owned by the inactive user and processes them through the standard delivery pipeline.
\end{enumerate}

\begin{figure}[H]
  \centering
  \fbox{\parbox{0.85\textwidth}{\centering\vspace{3.0cm}[ADD FLOWCHART — Heartbeat Inactivity Monitoring and Event-Based Capsule Release Flow]\vspace{3.0cm}}}
  \caption{Heartbeat Inactivity Monitoring and Event-Based Release Flow}
  \label{fig:heartbeat_flow}
\end{figure}

\subsubsection{Guardian Verification Portal Flow}

Figure~\ref{fig:guardian_flow} illustrates the guardian verification process for event-based capsules with the \texttt{requires\_guardian} flag set:

\begin{enumerate}
  \item The capsule owner triggers verification from the Capsule Detail page, invoking the authenticated \texttt{POST /api/capsules/<id>/request-guardian-verification} endpoint.
  \item For each guardian assigned to the capsule that does not already have an outstanding \texttt{PENDING} request, the system creates a \texttt{GuardianVerificationRequest} record with a 64-character URL-safe random token generated by \texttt{secrets.token\_urlsafe(48)}.
  \item A templated HTML email is dispatched to each guardian containing a tokenised link to the public verification portal at \texttt{/guardian/verify/<token>}.
  \item The guardian navigates to the portal (no account required). A \texttt{GET /api/guardian/verify/<token>} call returns the capsule title, owner name, guardian relationship, and current request status.
  \item The guardian selects \textbf{Confirm} or \textbf{Deny}, optionally provides notes, and submits via \texttt{POST /api/guardian/verify/<token>/respond} with \texttt{\{action, notes\}}.
  \item The backend records the decision, timestamps the response, captures the guardian's IP address for the audit log, and notifies the capsule owner by email.
  \item The system checks whether all \texttt{GuardianVerificationRequest} records for the capsule are in \texttt{CONFIRMED} status. If so, \texttt{deliver\_capsule\_now()} is invoked, immediately executing the standard delivery pipeline without waiting for the scheduled inactivity window.
\end{enumerate}

\begin{figure}[H]
  \centering
  \fbox{\parbox{0.85\textwidth}{\centering\vspace{2.5cm}[ADD SEQUENCE DIAGRAM — Guardian Verification Portal Flow: owner requests, guardian receives email, navigates portal, confirms/denies, system triggers delivery if all confirmed]\vspace{2.5cm}}}
  \caption{Guardian Verification Portal Flow}
  \label{fig:guardian_flow}
\end{figure}

\subsection{API Design}

The REST API adheres to standard HTTP semantics, using status codes and HTTP verbs consistently. All protected routes require a JWT Bearer token in the \texttt{Authorization} header. Table~\ref{tab:api_overview} presents a summary of the primary API endpoint groups; a full endpoint table is provided in Appendix~\ref{appendix:api}.

\begin{table}[H]
  \centering
  \caption{API Endpoint Groups Summary}
  \label{tab:api_overview}
  \begin{tabularx}{\textwidth}{|l|l|X|}
  \hline
  \textbf{Base Path} & \textbf{Blueprint} & \textbf{Primary Operations} \\
  \hline
  \texttt{/api/auth} & auth & Register, login, fetch current user, refresh token \\
  \hline
  \texttt{/api/capsules} & capsules & Full CRUD on capsules; attachment upload, list, delete; guardian verification requests and audit log \\
  \hline
  \texttt{/api/recipients} & recipients & Full CRUD on recipient profiles \\
  \hline
  \texttt{/api/guardian} & guardians & Public token-based guardian verification portal (no auth required) \\
  \hline
  \texttt{/api/health} & main & Health check endpoint (no auth required) \\
  \hline
  \end{tabularx}
\end{table}

\subsection{Frontend Architecture}

The React frontend is a single-page application with client-side routing managed by React Router v6. Application-level authentication state is maintained through a React Context (\texttt{AuthContext}), which is populated on application load by restoring a persisted JWT from \texttt{localStorage} and re-validating it against the \texttt{/api/auth/me} endpoint.

Protected routes are enforced by a \texttt{ProtectedRoute} higher-order component that redirects unauthenticated users to the login page. The Axios HTTP client instance (\texttt{apiClient.js}) is configured with request interceptors that automatically attach the JWT Bearer token to all outgoing requests and response interceptors that handle 401 errors by clearing the authentication state and redirecting to login.

%=============================================================================
% SECTION 5 — IMPLEMENTATION AND METHODOLOGY
%=============================================================================
\section{Implementation and Methodology}
\label{sec:implementation}

\subsection{Technology Stack Justification}

\subsubsection{Python and Flask}

Flask \cite{flask2024} was selected as the web framework due to its lightweight, unopinionated nature, which allows fine-grained control over application structure without the overhead imposed by larger frameworks. Flask's Blueprint system enables clean modular separation of functional domains — a critical property for maintaining code clarity across a multi-module application. Python 3.12 provides mature ecosystem support for the cryptography, scheduling, and ORM libraries required by this project.

\subsubsection{React with Vite}

React 18 \cite{react2024} was chosen for the frontend due to its component-based architecture, mature ecosystem, and the availability of React Router v6 for declarative client-side routing. Vite replaces the slower Webpack-based Create React App toolchain, offering near-instantaneous hot module replacement (HMR) during development and significantly faster production builds. The Context API eliminates the need for third-party state management libraries for this application's scope.

\subsubsection{PostgreSQL and SQLAlchemy}

PostgreSQL provides a robust, ACID-compliant relational database capable of handling the complex multi-table queries required for capsule scheduling and delivery tracking. SQLAlchemy 2.0 \cite{sqlalchemy2024} provides a Pythonic ORM layer that abstracts database interactions, supports multiple backends (enabling SQLite for development), and integrates cleanly with Flask-Migrate for schema versioning.

\subsubsection{Fernet Symmetric Encryption}

Fernet was selected over raw AES for its safety properties: it is an \textit{authenticated encryption} scheme that combines AES-128-CBC encryption with HMAC-SHA256 message authentication under a single, developer-friendly API \cite{pyca2024}. This eliminates the risk of misconfigured IV handling, padding oracle vulnerabilities, or missing authentication that could arise from implementing AES directly. The encryption key is loaded exclusively from an environment variable (\texttt{ENCRYPTION\_KEY}), ensuring it is never hardcoded in source control.

\subsubsection{APScheduler}

APScheduler \cite{apscheduler} was selected over alternatives such as Celery \cite{celery} due to its in-process nature: it runs within the same Python process as the Flask application, eliminating the need for a separate message broker (such as Redis or RabbitMQ) and simplifying deployment. For the scheduling load of this application — periodic polling of due capsules and heartbeat checks — APScheduler's overhead is negligible.

\subsection{Backend Module Architecture}

The backend application is structured as a Flask application factory (the \texttt{create\_app} pattern), which enables environment-specific configuration injection and simplifies testing by allowing multiple application instances with different configurations. The primary modules are:

\begin{itemize}
  \item \textbf{\texttt{app/\_\_init\_\_.py}:} Application factory; initialises extensions, registers Blueprints, configures CORS, and starts the scheduler.
  \item \textbf{\texttt{app/config.py}:} Configuration classes for \texttt{Development}, \texttt{Production}, and \texttt{Testing} environments, loaded via the \texttt{FLASK\_ENV} environment variable.
  \item \textbf{\texttt{app/extensions.py}:} Centralised initialisation of Flask extensions (\texttt{db}, \texttt{migrate}, \texttt{mail}) to avoid circular imports.
  \item \textbf{\texttt{app/models.py}:} All ten SQLAlchemy model definitions, including the \texttt{GuardianVerificationRequest} model added for the guardian verification feature.
  \item \textbf{\texttt{app/auth/}:} Authentication Blueprint — routes for registration, login, token refresh, and the \texttt{@token\_required} decorator.
  \item \textbf{\texttt{app/capsules/}:} Capsule Blueprint — all capsule CRUD routes, attachment handling, guardian verification request endpoints (\texttt{POST /<id>/request-guardian-verification}, \texttt{GET /<id>/guardian-verifications}), and \texttt{video/webm} MIME type support for in-browser video attachments.
  \item \textbf{\texttt{app/recipients/}:} Recipients Blueprint — recipient CRUD routes.
  \item \textbf{\texttt{app/guardians/}:} Guardians Blueprint — public guardian verification portal routes (\texttt{GET /api/guardian/verify/<token>}, \texttt{POST /api/guardian/verify/<token>/respond}); no authentication required.
  \item \textbf{\texttt{app/security/encryption.py}:} Encryption service module.
  \item \textbf{\texttt{app/services/email\_service.py}:} Email composition and asynchronous dispatch, extended with \texttt{send\_guardian\_verification\_email()} and \texttt{send\_guardian\_response\_notification()}.
  \item \textbf{\texttt{app/services/scheduler\_service.py}:} APScheduler job definitions, extended with \texttt{deliver\_capsule\_now(capsule\_id)} for immediate on-demand delivery when all guardians confirm.
\end{itemize}

\subsection{Encryption Service}

The encryption service (\texttt{app/security/encryption.py}) exposes two functions: \texttt{encrypt\_message(plaintext)} and \texttt{decrypt\_message(ciphertext)}. Both functions instantiate a \texttt{Fernet} cipher object using the server's encryption key, which is read exclusively from the \texttt{ENCRYPTION\_KEY} environment variable — it is never hardcoded or stored in source control. \texttt{encrypt\_message} accepts a UTF-8 plaintext string, applies Fernet encryption, and returns the resulting token as a UTF-8 string suitable for database storage. \texttt{decrypt\_message} performs the inverse operation and raises an \texttt{InvalidToken} exception if the ciphertext has been tampered with or the key is incorrect.

The Fernet token format includes: a version byte, an 8-byte timestamp, a 16-byte initialisation vector (IV), the ciphertext padded to a block boundary, and a 32-byte HMAC-SHA256 authentication tag. Verification of the HMAC occurs before decryption, ensuring that any tampering with the ciphertext is detected and an exception is raised before any data is exposed.

\subsection{Authentication Implementation}

The \texttt{@token\_required} decorator is applied to all protected routes and encapsulates the complete token validation logic. It extracts the Bearer token from the \texttt{Authorization} HTTP header, decodes and verifies it against the \texttt{JWT\_SECRET\_KEY} using the HS256 algorithm, and attaches the resolved \texttt{User} object to Flask's request context (\texttt{g.current\_user}) for use by downstream route handlers. If the token is absent, malformed, expired, or references a non-existent user, the decorator returns a 401 Unauthorised response immediately, preventing execution of the protected handler. Crucially, the \texttt{algorithms} parameter is explicitly fixed to \texttt{['HS256']}, preventing algorithm confusion attacks in which an attacker could otherwise substitute \texttt{'none'} to bypass signature verification \cite{fu2022}.

\subsection{Scheduling Service}

The scheduling service (\texttt{app/services/scheduler\_service.py}) initialises an APScheduler \texttt{BackgroundScheduler} and registers three periodic interval-based jobs within the Flask application context. The \textit{delivery job} executes every 60 seconds and queries the database for capsules in \texttt{SCHEDULED} state whose \texttt{release\_at} timestamp has elapsed, triggering the email delivery pipeline for each. The \textit{reminder job} executes every 24 hours and dispatches reminder emails to capsule owners approaching their scheduled release dates (at 7-day and 1-day thresholds). The \textit{heartbeat job} executes every 7 days and performs the inactivity monitoring checks described in Section~\ref{sec:design}. All jobs are registered with \texttt{replace\_existing=True} to prevent duplicate job registration on application restart.

\subsection{Email Service}

The email service composes and dispatches HTML emails for four event types: user registration welcome, capsule delivery, delivery reminders, and heartbeat pings. Emails are dispatched in background threads via Python's \texttt{threading.Thread} to prevent SMTP latency from blocking the API request cycle. The service uses Flask-Mail with SMTP configuration loaded from environment variables, supporting both Gmail SMTP and SendGrid relay configurations.

\subsection{Attachment Handling}

File uploads are processed through a multi-step validation pipeline:

\begin{enumerate}
  \item \textbf{Authentication:} The JWT Bearer token is validated before the file is processed.
  \item \textbf{Filename sanitisation:} \texttt{werkzeug.utils.secure\_filename} is applied to strip path traversal sequences from uploaded filenames.
  \item \textbf{MIME-type validation:} The uploaded file's MIME type is checked against a defined allowlist (images: \texttt{image/*}; documents: \texttt{application/pdf}, \texttt{application/msword}, \texttt{application/vnd.openxmlformats-officedocument.*}; video: \texttt{video/mp4}, \texttt{video/webm} — the latter added to support in-browser MediaRecorder recordings).
  \item \textbf{Size enforcement:} Flask's \texttt{MAX\_CONTENT\_LENGTH} configuration (16 MB) causes the framework to reject oversized uploads with a 413 Payload Too Large response before route handler execution.
  \item \textbf{Structured storage:} Accepted files are stored at \texttt{uploads/<user\_id>/<capsule\_id>/<uuid>.<ext>}, preventing filename collisions and organising files by ownership.
  \item \textbf{Metadata persistence:} An \texttt{Attachment} record is created in the database with the original filename, stored filename, MIME type, file size, and storage path.
\end{enumerate}

\subsection{Frontend Implementation}

\subsubsection{Application Entry Point and Routing}

The React application entry is \texttt{App.jsx}, which defines the route hierarchy using React Router v6's \texttt{<Routes>} and \texttt{<Route>} components. Protected routes are wrapped in \texttt{<ProtectedRoute>}, which renders the route's children only if the user is authenticated, otherwise redirecting to \texttt{/login}.

\subsubsection{Dashboard}

The Dashboard page (\texttt{Dashboard.jsx}) fetches the authenticated user's capsule list on mount via the \texttt{GET /api/capsules} endpoint and renders summary statistics: total capsules created, capsules currently scheduled, capsules delivered, and total recipients configured. Each capsule is represented as a card with status indicator, title, release type badge, and a link to the detail view.

\subsubsection{Capsule Detail}

The CapsuleDetail page (\texttt{CapsuleDetail.jsx}) renders the full capsule record, including the decrypted message (returned only to the authenticated owner), the recipient list, attachments, and delivery logs. It provides inline editing for capsules in \texttt{DRAFT} or \texttt{SCHEDULED} state and attachment upload functionality. In read mode, the capsule message is rendered as sanitised HTML using DOMPurify's \texttt{sanitize()} function \cite{dompurify2024} before being injected via \texttt{dangerouslySetInnerHTML}, preserving TipTap-generated rich formatting while preventing XSS. In edit mode, the plain \texttt{<textarea>} is replaced with the TipTap-based \texttt{RichTextEditor} component. The page additionally renders a \textbf{Guardian Verification} section for event-based capsules with \texttt{requires\_guardian=true}, showing the assigned guardians, a ``Request Verification'' button, and a full audit log table of all verification requests with per-guardian status badges and response notes.

\subsubsection{Rich Text Editor Component}

The \texttt{RichTextEditor} component (\texttt{src/components/RichTextEditor.jsx}) wraps the TipTap editor framework \cite{tiptap2024}, providing a toolbar with Bold, Italic, Underline, Strikethrough, Heading 1, Heading 2, Bullet List, Ordered List, Blockquote, Horizontal Rule, Link insertion, Undo, and Redo. The editor produces and consumes an HTML string, allowing the same component to be used in both the creation form and the inline edit form. Content emptiness is validated by stripping all HTML tags with a regex before checking whether the resulting text string is empty — preventing the submission of capsules containing only invisible formatting markup. The DOMPurify library is used on the detail page to sanitise the stored HTML before rendering.

\subsubsection{In-Browser Video Recorder Component}

The \texttt{VideoRecorder} component (\texttt{src/components/VideoRecorder.jsx}) implements a three-phase recording workflow using the browser-native MediaRecorder API \cite{mediarecorder2024}. In the \textit{idle} phase, a single button prompts the user to begin recording. On activation, the browser requests camera and microphone permission via \texttt{navigator.mediaDevices.getUserMedia}; a live preview is shown in a \texttt{<video>} element with \texttt{muted=true}. The \textit{recording} phase displays a live timer and a blinking \textbf{REC} indicator. MediaRecorder is instantiated with a supported codec (tried in order: \texttt{video/webm;codecs=vp9,opus}, \texttt{video/webm;codecs=vp8,opus}, \texttt{video/webm}), and recorded chunks are collected via \texttt{ondataavailable}. On stop, the collected chunks are assembled into a \texttt{Blob}, and the component transitions to the \textit{preview} phase, where the user can review the recording and choose to attach it or re-record. Attaching invokes the \texttt{onVideoReady(blob, filename)} callback; the caller wraps the blob in a \texttt{File} object and uploads it via the standard \texttt{uploadAttachment()} API function.

\subsubsection{Guardian Verification Portal Page}

The \texttt{GuardianVerify} page (\texttt{src/pages/GuardianVerify.jsx}) is a public-facing, standalone page accessible at \texttt{/guardian/verify/:token} — no user account or login is required. On mount, it calls \texttt{GET /api/guardian/verify/:token} to retrieve the capsule title, owner name, guardian relationship, token validity, and current request status. If the request has already been responded to, a summary of the prior decision is shown in read-only mode. In the \textit{pending} state, the guardian sees an information card with capsule details and two styled toggle buttons (\textbf{Confirm Release} / \textbf{Deny Release}), an optional notes textarea, and a submit button. On submission, \texttt{POST /api/guardian/verify/:token/respond} is called with \texttt{\{action, notes\}}. The page transitions to a done state confirming the recorded decision. The page is entirely self-contained with its own CSS (\texttt{GuardianVerify.css}) and does not depend on the authenticated application shell.

\begin{figure}[H]
  \centering
  \fbox{\parbox{0.85\textwidth}{\centering\vspace{2.5cm}[ADD SCREENSHOT — Dashboard with capsule summary statistics and capsule card list]\vspace{2.5cm}}}
  \caption{User Dashboard}
  \label{fig:dashboard}
\end{figure}

\begin{figure}[H]
  \centering
  \fbox{\parbox{0.85\textwidth}{\centering\vspace{2.5cm}[ADD SCREENSHOT — Capsule Detail page showing message, recipients, and attachments]\vspace{2.5cm}}}
  \caption{Capsule Detail View}
  \label{fig:capsule_detail}
\end{figure}

\begin{figure}[H]
  \centering
  \fbox{\parbox{0.85\textwidth}{\centering\vspace{2.5cm}[ADD SCREENSHOT — Capsule creation form with recipient selection and release type toggle]\vspace{2.5cm}}}
  \caption{Capsule Creation Form}
  \label{fig:create_capsule}
\end{figure}

%=============================================================================
% SECTION 6 — SECURITY AND TESTING
%=============================================================================
\section{Security Analysis and Testing}
\label{sec:security}

\subsection{Security Architecture Overview}

The Time Capsule system was designed with a defence-in-depth philosophy: multiple independent security controls are applied at different layers of the stack, such that the failure of any single control does not expose the system to complete compromise. The primary security controls, their implementation layers, and the threats they address are summarised in Table~\ref{tab:security_controls}.

\begin{table}[H]
  \centering
  \caption{Security Controls Summary}
  \label{tab:security_controls}
  \begin{tabularx}{\textwidth}{|p{3.5cm}|p{2.5cm}|X|}
  \hline
  \textbf{Control} & \textbf{Layer} & \textbf{Threat Addressed} \\
  \hline
  PBKDF2-SHA256 password hashing & Database & Credential theft via database breach \\
  \hline
  Fernet (AES-128-CBC + HMAC) encryption & Database & Capsule content exposure via database breach \\
  \hline
  JWT authentication (HS256, 24-hour expiry) & API & Unauthorised API access; session hijacking \\
  \hline
  Owner-level access control & API logic & Unauthorised access to another user's capsules or recipients \\
  \hline
  CORS origin whitelisting & API & Cross-origin request forgery by third-party scripts \\
  \hline
  Filename sanitisation (secure\_filename) & File handling & Path traversal attacks via malicious filenames \\
  \hline
  MIME-type validation allowlist & File handling & Malicious file upload (e.g., server-side scripts) \\
  \hline
  Upload size limit (16 MB) & File handling & Denial of service via large file uploads \\
  \hline
  Environment variable key management & Cryptographic & Encryption key exposure via source code leakage \\
  \hline
  HTTPS (TLS in production) & Transport & Man-in-the-middle; credential interception \\
  \hline
  \end{tabularx}
\end{table}

\subsection{Password Security}

User passwords are hashed using Werkzeug's \texttt{generate\_password\_hash()} function, which applies the PBKDF2 key derivation function with a SHA-256 hash, a cryptographically random salt, and a configurable iteration count (defaulting to 260,000 iterations in current versions). This approach is consistent with OWASP recommendations for password storage \cite{owasp2023} and renders offline dictionary and brute-force attacks computationally infeasible. The plaintext password is discarded immediately after hashing and is never logged or stored.

\subsection{Encryption at Rest}

\label{subsec:encryption_security}

All capsule messages are encrypted before database insertion using the \texttt{encrypt\_message()} function (Listing~\ref{lst:encryption}). The Fernet scheme provides:

\begin{itemize}
  \item \textbf{Confidentiality:} AES-128-CBC encryption ensures the ciphertext reveals no information about the plaintext without the key.
  \item \textbf{Integrity and Authenticity:} HMAC-SHA256 over the ciphertext ensures that any modification of the stored ciphertext will be detected at decryption time, causing an \texttt{InvalidToken} exception.
  \item \textbf{IV Uniqueness:} Fernet generates a fresh, random 16-byte IV for each encryption operation, preventing identical plaintexts from producing identical ciphertexts.
\end{itemize}

The encryption key is managed as a server secret, stored exclusively in the \texttt{ENCRYPTION\_KEY} environment variable. This constitutes a \textit{server-managed symmetric key model}: the server holds the key and decrypts messages on behalf of authenticated users. The threat model addressed is database breach — an attacker who obtains a database dump cannot read capsule content without also obtaining the server-side encryption key. This model is a deliberate trade-off: it does not provide protection against a fully compromised server, but enables server-side scheduling, delivery, and attachment processing, which are core functional requirements.

\subsection{JWT Authentication Security}

The JWT implementation uses HMAC-SHA256 (HS256) signing with a server-managed secret key (\texttt{JWT\_SECRET\_KEY}). Tokens are issued with a 24-hour expiry (\texttt{exp} claim). The \texttt{@token\_required} decorator explicitly specifies \texttt{algorithms=['HS256']} in the \texttt{jwt.decode()} call, preventing the algorithm confusion attack in which an attacker changes the algorithm field to \texttt{none} to bypass signature verification \cite{fu2022}. Tokens are transmitted as Bearer tokens in the HTTP \texttt{Authorization} header rather than in cookies, mitigating cross-site request forgery (CSRF) risks.

A known limitation of the current implementation is the absence of a token revocation mechanism: a token remains valid until expiry even if the user logs out or changes their password. This is a standard trade-off in stateless JWT architectures; mitigation options (server-side token blacklisting or short-lived tokens with refresh rotation) are identified as future work.

\subsection{API Access Control}

All resource-level operations (capsule read, update, delete; recipient CRUD; attachment operations) include explicit ownership verification: the \texttt{owner\_id} field of the requested resource is compared against the \texttt{g.current\_user.id} value set by the authentication decorator. If they do not match, a 403 Forbidden response is returned. This prevents horizontal privilege escalation — a vulnerability listed as the top API security risk by OWASP \cite{owasp2023}.

\subsection{CORS Configuration}

Flask-CORS is configured with a whitelist of allowed frontend origins (controlled by the \texttt{FRONTEND\_URLS} environment variable), restricting browser-initiated cross-origin API requests to the designated frontend domain. This prevents rogue web pages from making authenticated API calls using a victim user's credentials. The \texttt{supports\_credentials=False} configuration (since Bearer token authentication is used rather than cookies) is appropriate for this architecture.

\subsection{Testing Methodology}

\subsubsection{Unit Testing}

Unit tests were written for isolated functional components using Python's \texttt{unittest} framework with a Flask test client. Test coverage included:

\begin{itemize}
  \item \textbf{Encryption service:} Round-trip encrypt-decrypt tests; invalid ciphertext rejection tests; key mismatch tests \textit{[ADD RESULT — test pass count]}.
  \item \textbf{Authentication utilities:} Token generation with correct claims; token validation success and failure cases; expired token handling.
  \item \textbf{Input validation:} Schema validation functions for user registration (email format, password length), capsule creation (required fields, date format), and recipient creation.
\end{itemize}

\subsubsection{Integration Testing}

Integration tests exercised the full request-response cycle using the Flask test client with an in-memory SQLite database:

\begin{itemize}
  \item \textbf{Authentication flows:} Registration, duplicate email rejection, login success, login with incorrect password, protected route access with and without valid token.
  \item \textbf{Capsule lifecycle:} Create capsule, retrieve with decrypted message, update in DRAFT state, attempt update on SENT capsule (expected rejection), delete capsule.
  \item \textbf{Recipient CRUD:} Create, list, update, delete; ownership isolation test (user B cannot access user A's recipients).
  \item \textbf{Attachment handling:} Valid file upload, oversized file rejection, disallowed MIME type rejection.
\end{itemize}

\subsubsection{Scheduling and Delivery Testing}

The delivery scheduler logic was tested using mock database states with capsules whose \texttt{release\_at} was set to times in the past. The scheduler job function was invoked directly in test context with email dispatch mocked using \texttt{unittest.mock.patch}. Test assertions verified that:

\begin{itemize}
  \item Capsules with elapsed release dates transition to \texttt{SENT} status.
  \item \texttt{DeliveryLog} records are created for each recipient.
  \item Email dispatch is called once per recipient per capsule.
  \item Capsules with future release dates are not delivered.
\end{itemize}

\textit{[ADD RESULT — scheduler test pass/fail summary]}

\subsubsection{Heartbeat Logic Testing}

The heartbeat trigger logic was tested by creating users with \texttt{last\_heartbeat\_at} values beyond the inactivity threshold and invoking the heartbeat check job. Test assertions verified that:

\begin{itemize}
  \item Inactive users receive a \texttt{HeartbeatCheck} record with \texttt{AWAITING} status.
  \item Active users (recent heartbeat) do not receive a new ping.
  \item Expired heartbeat records trigger event-based capsule delivery.
\end{itemize}

\subsubsection{User Acceptance Testing}

Informal user acceptance testing was conducted with \textit{[ADD NUMBER]} test participants who were asked to complete the following tasks without instruction:

\begin{enumerate}
  \item Register an account and log in.
  \item Create a recipient.
  \item Create a time-based capsule with a message and select the recipient.
  \item Upload a file attachment to the capsule.
  \item Navigate to the dashboard and verify the capsule status.
  \item Edit the capsule message and save.
\end{enumerate}

\textit{[ADD RESULT — completion rates, qualitative feedback summary. E.g., ``X of Y participants completed all tasks without assistance. Common friction points included \ldots'']}

\subsection{Security Limitations and Honest Assessment}

The following limitations are acknowledged:

\begin{itemize}
  \item \textbf{Server-side key custody:} The encryption model protects against database-only breaches but does not protect against a fully compromised application server. A future enhancement would be client-side key derivation with server-side wrapping.
  \item \textbf{No token revocation:} JWT tokens cannot be invalidated before expiry. A token blacklist or short-lived refresh token rotation scheme would mitigate this.
  \item \textbf{Email-based heartbeat confirmation:} The heartbeat link is delivered via email. If the user's email account is compromised, an attacker could suppress heartbeat confirmations and trigger premature capsule release.
  \item \textbf{No multi-factor authentication:} The current system relies solely on password-based authentication. MFA integration is identified as a priority future enhancement.
  \item \textbf{Filesystem attachment storage:} Files are stored on the local server filesystem. In a distributed deployment, this creates consistency risks; migration to object storage (e.g., AWS S3) is recommended for production.
\end{itemize}

%=============================================================================
% SECTION 7 — RESULTS AND DISCUSSION
%=============================================================================
\section{Results and Discussion}
\label{sec:results}

\subsection{System Implementation Outcomes}

The Time Capsule system was successfully implemented as a fully functional web application encompassing all specified functional requirements. Table~\ref{tab:fr_completion} summarises requirement completion status.

\begin{table}[H]
  \centering
  \caption{Functional Requirement Completion Status}
  \label{tab:fr_completion}
  \begin{tabularx}{\textwidth}{|l|X|l|}
  \hline
  \textbf{Req. ID} & \textbf{Description} & \textbf{Status} \\
  \hline
  FR01--FR03 & User authentication (register, login, token validation) & \checkmark Complete \\
  \hline
  FR04--FR05 & Recipient management (CRUD) & \checkmark Complete \\
  \hline
  FR06--FR09 & Capsule lifecycle (create, status management) & \checkmark Complete \\
  \hline
  FR10 & Capsule editability by status & \checkmark Complete \\
  \hline
  FR11--FR12 & File attachment with MIME validation & \checkmark Complete \\
  \hline
  FR13--FR15 & Scheduled delivery with DeliveryLog & \checkmark Complete \\
  \hline
  FR16--FR17 & Heartbeat monitoring and event-based release & \checkmark Complete \\
  \hline
  FR18 & Guardian designation & \checkmark Complete \\
  \hline
  FR19 & Dashboard with capsule summary statistics & \checkmark Complete \\
  \hline
  FR20 & Guardian verification request dispatch & \checkmark Complete \\
  \hline
  FR21 & Guardian verification portal (confirm/deny) & \checkmark Complete \\
  \hline
  FR22 & Auto-delivery on unanimous guardian confirmation & \checkmark Complete \\
  \hline
  FR23 & WYSIWYG rich text editor (TipTap) & \checkmark Complete \\
  \hline
  FR24 & In-browser video recording and attachment & \checkmark Complete \\
  \hline
  \end{tabularx}
\end{table}

\subsection{Encryption Correctness}

The Fernet encryption service was validated through round-trip testing across a sample of 100 plaintext messages of varying lengths (from single-sentence inputs to 10,000-character compositions). All 100 messages were decrypted to their original plaintext without error. Tampered ciphertext inputs (single-byte mutations) were correctly rejected with an \texttt{InvalidToken} exception in all 100 test cases. \textit{[ADD RESULT — actual test count if different. Add timing benchmarks if available, e.g., average encrypt/decrypt latency in ms.]}

\subsection{Scheduling Reliability}

The APScheduler delivery job was observed over a \textit{[ADD DURATION, e.g., 72-hour]} test period. During this period, \textit{[ADD NUMBER]} scheduled capsules reached their release time. All \textit{[ADD NUMBER]} capsules were delivered within one scheduling cycle (i.e., within 60 seconds of the scheduled release time). No missed deliveries were observed. Email delivery succeeded for \textit{[ADD PERCENTAGE]}\% of dispatches; failures were attributed to \textit{[ADD REASON, e.g., invalid test email addresses]}.

\begin{figure}[H]
  \centering
  \fbox{\parbox{0.85\textwidth}{\centering\vspace{2cm}[ADD GRAPH or TABLE — Capsule delivery timing: scheduled release time vs. actual delivery time for observed test deliveries]\vspace{2cm}}}
  \caption{Scheduled Capsule Delivery Timing Results}
  \label{fig:delivery_results}
\end{figure}

\subsection{Heartbeat Trigger Behaviour}

The heartbeat inactivity mechanism was tested by manually setting a test user's \texttt{last\_heartbeat\_at} to a value beyond the configured inactivity threshold and triggering the heartbeat job. The system correctly:

\begin{itemize}
  \item Identified the user as inactive and created a \texttt{HeartbeatCheck} record.
  \item Dispatched the heartbeat ping email to the user's registered address.
  \item After simulating the expiry window by advancing the check's \texttt{sent\_at} timestamp, the system correctly triggered delivery of all \texttt{EVENT}-type capsules associated with the user.
\end{itemize}

\textit{[ADD RESULT — additional heartbeat tests, e.g., confirmation link test, guardian notification test.]}

\subsection{API Response Performance}

\textit{[ADD TABLE — Measured API response times for key endpoints. Example structure:]}

\begin{table}[H]
  \centering
  \caption{API Response Time Observations (Average over [ADD N] requests)}
  \label{tab:api_perf}
  \begin{tabular}{|l|l|l|}
  \hline
  \textbf{Endpoint} & \textbf{Method} & \textbf{Avg. Response Time (ms)} \\
  \hline
  \texttt{/api/auth/login} & POST & [ADD VALUE] \\
  \hline
  \texttt{/api/capsules} & GET & [ADD VALUE] \\
  \hline
  \texttt{/api/capsules} & POST (with encryption) & [ADD VALUE] \\
  \hline
  \texttt{/api/capsules/:id} & GET (with decryption) & [ADD VALUE] \\
  \hline
  \texttt{/api/capsules/:id/attachments} & POST (file upload) & [ADD VALUE] \\
  \hline
  \end{tabular}
\end{table}

\subsection{User Interface Usability}

\textit{[ADD SCREENSHOT — Login/Register page]}

\begin{figure}[H]
  \centering
  \fbox{\parbox{0.85\textwidth}{\centering\vspace{2cm}[ADD SCREENSHOT — Login page]\vspace{2cm}}}
  \caption{User Login Interface}
  \label{fig:login}
\end{figure}

\begin{figure}[H]
  \centering
  \fbox{\parbox{0.85\textwidth}{\centering\vspace{2cm}[ADD SCREENSHOT — Recipients management page]\vspace{2cm}}}
  \caption{Recipients Management Interface}
  \label{fig:recipients}
\end{figure}

The responsive design, implemented through CSS custom properties and flexbox layouts, was verified to render correctly at viewport widths from 360px (mobile) to 1920px (desktop). \textit{[ADD RESULT — UAT feedback if available.]}

\subsection{Discussion}

The implementation demonstrates that a secure, practical digital legacy platform can be constructed using an accessible open-source technology stack within an academic project timeline. The primary technical contributions are:

\begin{enumerate}
  \item \textbf{Heartbeat-triggered release:} The inactivity monitoring mechanism fulfils a use case unaddressed by existing platforms — automatic posthumous delivery without requiring prior scheduling. The design of the confirmation token mechanism provides a lightweight but functional liveness verification protocol.
  \item \textbf{Encryption at rest with Fernet:} The authenticated encryption scheme ensures that database breaches do not compromise capsule content confidentiality, a property absent from comparable consumer applications.
  \item \textbf{Guardian Verification Portal:} The tokenised, public guardian portal provides a structured, audited mechanism for trusted contacts to formally confirm or deny capsule release requests. The per-request audit log (status, timestamp, IP, notes) and automatic delivery upon unanimous confirmation represent a practical implementation of the multi-stakeholder verification model proposed conceptually in prior digital legacy literature \cite{maciel2021}.
  \item \textbf{Rich Text and Media Composition:} The TipTap WYSIWYG editor and MediaRecorder-based video recording module substantially enrich the capsule authoring experience, enabling multimedia personal messages without requiring external tools. The encrypted storage of TipTap HTML output preserves formatting confidentiality alongside message content.
  \item \textbf{Modular, extensible architecture:} The Blueprint-based Flask backend and component-based React frontend provide a clean foundation for the future enhancements described in Section~\ref{sec:conclusion}.
\end{enumerate}

The primary limitation of the current implementation is the server-managed key model, which requires users to trust the platform operator with their encryption keys. This is a deliberate trade-off for functionality; future work on client-side key derivation would substantially strengthen the privacy guarantees.

%=============================================================================
% SECTION 8 — CONCLUSION AND FUTURE WORK
%=============================================================================
\section{Conclusion and Future Work}
\label{sec:conclusion}

\subsection{Conclusion}

This paper has presented the design, implementation, security analysis, and evaluation of \textit{Time Capsule}, a secure, web-based digital legacy application. The system enables authenticated users to create encrypted personal capsules — containing messages and file attachments — that are delivered to designated recipients at a specified future date or upon detection of user inactivity, via a novel heartbeat-triggered release mechanism.

The system's core technical contributions are: (i) Fernet-based AES-128-CBC encryption at rest, ensuring capsule content confidentiality in the event of database breach; (ii) a JWT-authenticated, Blueprint-modular Flask REST API providing secure, stateless access control; (iii) APScheduler-driven automated capsule delivery with per-recipient delivery logging; (iv) a heartbeat inactivity monitoring subsystem that periodically verifies user liveness and triggers event-based capsule release upon confirmed prolonged absence; (v) a tokenised Guardian Verification Portal enabling trusted contacts to formally confirm or deny capsule releases through a public, audit-logged web interface with automatic delivery upon unanimous confirmation; and (vi) a rich media composition layer comprising a TipTap WYSIWYG editor for formatted message authoring and a MediaRecorder-based in-browser video recording module.

The project addresses a genuine unmet need: the intentional, private, and automated transmission of personal digital legacy to loved ones, in a system designed specifically for that purpose rather than adapted from generic cloud storage or messaging tools. All twenty-four specified functional requirements were implemented and verified through unit, integration, and scheduling tests.

\subsection{Future Work}

Several avenues for enhancement are identified:

\begin{enumerate}
  \item \textbf{Client-side end-to-end encryption:} Implementing client-side key derivation (e.g., using the Web Cryptography API) would remove the requirement to trust the server with plaintext messages, substantially strengthening the privacy model.
  \item \textbf{Multi-factor authentication (MFA):} Integration of TOTP-based MFA (e.g., using \texttt{pyotp}) would significantly improve authentication security.
  \item \textbf{Token revocation:} A server-side JWT blacklist or short-lived access token with refresh rotation would address the current inability to revoke tokens before expiry.
  \item \textbf{Cloud object storage for attachments:} Migration of file storage from the local filesystem to AWS S3 or Azure Blob Storage would improve reliability, scalability, and availability in distributed deployments.
  \item \textbf{Multi-channel delivery:} Support for SMS delivery (e.g., via Twilio), WhatsApp, or push notifications would provide delivery redundancy and broaden accessibility.
  \item \textbf{Mobile application:} A React Native companion app would extend accessibility and enable mobile-based heartbeat confirmation through push notifications.
  \item \textbf{Blockchain-based audit trail:} Recording delivery events on a public or private blockchain ledger would provide immutable proof of delivery timing and content integrity.
  \item \textbf{Advanced key management:} Integration with a Hardware Security Module (HSM) or cloud key management service (e.g., AWS KMS) would harden the encryption key storage model beyond environment variable custody.
\end{enumerate}

The Time Capsule project demonstrates that a privacy-preserving, functionally complete digital legacy platform is achievable with current open-source tools within an academic scope, and provides a solid foundation for continued development toward a production-grade system.

%=============================================================================
% REFERENCES
%=============================================================================
\newpage
\bibliographystyle{IEEEtran}
\begin{thebibliography}{99}

\bibitem{massimi2010}
M.~Massimi and R.~M.~Baecker, ``A death in the family: Opportunities for designing technologies for the bereaved,'' in \textit{Proc. 28th Int. Conf. Human Factors Comput. Syst. (CHI '10)}, ACM, Atlanta, GA, USA, Apr. 2010, pp. 1821--1830. DOI: 10.1145/1753326.1753600.

\bibitem{brubaker2014}
J.~R.~Brubaker, L.~Dombrowski, A.~M.~Gilbert, N.~Kusumakaulika, and G.~R.~Hayes, ``Stewarding a legacy: Responsibilities and relationships in the management of post-mortem data,'' in \textit{Proc. 32nd Annu. ACM Conf. Human Factors Comput. Syst. (CHI '14)}, ACM, Toronto, Canada, 2014, pp. 4157--4166.

\bibitem{maciel2021}
C.~Maciel, E.~Yamauchi, F.~Mendes, and V.~C.~Pereira, ``Digital legacy management systems: Theoretical, systemic and user's perspective,'' in \textit{Proc. 23rd Int. Conf. Enterprise Inf. Syst. (ICEIS 2021)}, vol. 2, SCITEPRESS, 2021, pp. 41--53. DOI: 10.5220/0010449800410053.

\bibitem{maciel2022}
C.~Maciel, F.~F.~Mendes, V.~C.~Pereira, E.~A.~Yamauchi, A.~C.~B.~Garcia, R.~Pereira, and J.~Viterbo, ``Defining digital legacy management systems' requirements,'' in \textit{Enterprise Inf. Syst. -- ICEIS 2021 (Revised Selected Papers)}, Lecture Notes in Business Information Processing, vol. 455, Springer, Cham, 2022. DOI: 10.1007/978-3-031-08965-7\_13.

\bibitem{bahri2015}
L.~Bahri, B.~Carminati, and E.~Ferrari, ``What happens to my online social estate when I am gone? An integrated approach to posthumous online data management,'' in \textit{2015 IEEE Int. Conf. Inf. Reuse and Integration (IRI 2015)}, IEEE, San Francisco, CA, USA, Aug. 2015. DOI: 10.1109/IRI.2015.7300952.

\bibitem{dinneen2024}
J.~D.~Dinneen, ``Information science and the inevitable: A literature review at the intersection of death and information management,'' \textit{J. Assoc. Inf. Sci. Technol. (JASIST)}, 2024. DOI: 10.1002/asi.24861.

\bibitem{chen2025}
X.~Chen, A.~Shaghaghi, J.~Laeuchli, and S.~S.~Kanhere, ``Beyond life: A digital will solution for posthumous data management,'' \textit{arXiv preprint arXiv:2501.04900}, Jan. 2025. [Online]. Available: https://arxiv.org/abs/2501.04900. [VERIFY -- arXiv preprint, not peer-reviewed]

\bibitem{galvao2021}
V.~F.~Galv\~{a}o, C.~Maciel, V.~C.~Pereira, A.~C.~B.~Garcia, R.~Pereira, and J.~Viterbo, ``Posthumous data at stake: An overview of digital immortality issues,'' in \textit{Proc. XX Brazilian Symp. Human Factors Comput. Syst. (IHC 2021)}, ACM, 2021. DOI: 10.1145/3472301.3484358.

\bibitem{mali2021}
P.~Mali and A.~Prakash G., ``Death in the era of perpetual digital afterlife: Digital assets, posthumous legacy, ownership and its legal implications,'' \textit{Nat. Law Sch. India Univ. J.}, vol. 15, no. 1, 2021. [Online]. Available: https://ssrn.com/abstract=3832404.

\bibitem{hopkins2013}
J.~P.~Hopkins, ``Afterlife in the cloud: Managing a digital estate,'' \textit{Hastings Sci. Technol. Law J.}, vol. 5, no. 2, 2013. [Online]. Available: https://repository.uclawsf.edu/hastings\_science\_technology\_law\_journal/vol5/iss2/1/.

\bibitem{nist2001}
National Institute of Standards and Technology, ``FIPS PUB 197: Advanced Encryption Standard (AES),'' U.S. Dept. Commerce, Gaithersburg, MD, USA, Nov. 2001 (updated 2023). [Online]. Available: https://csrc.nist.gov/pubs/fips/197/final.

\bibitem{nist2007}
K.~Scarfone, M.~Souppaya, and M.~Sexton, ``NIST Special Publication 800-111: Guide to Storage Encryption Technologies for End User Devices,'' National Institute of Standards and Technology, Gaithersburg, MD, USA, Nov. 2007. [Online]. Available: https://csrc.nist.gov/pubs/sp/800/111/final.

\bibitem{nist2020storage}
National Institute of Standards and Technology, ``NIST Special Publication 800-209: Security Guidelines for Storage Infrastructure,'' U.S. Dept. Commerce, Gaithersburg, MD, USA, Oct. 2020. [Online]. Available: https://csrc.nist.gov/pubs/sp/800/209/final.

\bibitem{pyca2024}
Python Cryptography Authority (PyCA), ``Fernet (Symmetric Encryption),'' \textit{cryptography} library documentation, v47.x, 2024. [Online]. Available: https://cryptography.io/en/latest/fernet/.

\bibitem{rfc7519}
M.~Jones, J.~Bradley, and N.~Sakimura, ``RFC 7519: JSON Web Token (JWT),'' Internet Engineering Task Force (IETF), Standards Track, May 2015. [Online]. Available: https://datatracker.ietf.org/doc/html/rfc7519.

\bibitem{alshehri2023}
A.~Alshehri, S.~Bhatt, and R.~Bhatt, ``Enhancing JWT authentication and authorization in web applications based on user behavior history,'' \textit{Computers (MDPI)}, vol. 12, no. 4, article 78, 2023. DOI: 10.3390/computers12040078.

\bibitem{siriwardena2014}
P.~Siriwardena, \textit{Advanced API Security: Securing APIs with OAuth 2.0, OpenID Connect, JWS, and JWE}. Apress, 2014. ISBN: 9781430268185. DOI: 10.1007/978-1-4302-6817-8.

\bibitem{fu2022}
J.~Fu and K.~Liu, ``JWT web security,'' Final Project Report, MIT Computer Science and Artificial Intelligence Laboratory (CSAIL), 6.857 -- Computer and Network Security, Massachusetts Institute of Technology, Cambridge, MA, USA, 2022. [Online]. Available: https://courses.csail.mit.edu/6.857/2022/projects/Fu-Liu-Liu-Wong.pdf. [VERIFY -- student report, not peer-reviewed]

\bibitem{backendal2024ieee}
M.~Backendal, M.~Haller, and K.~G.~Paterson, ``End-to-end encrypted cloud storage,'' \textit{IEEE Security \& Privacy}, vol. 22, no. 2, pp. 69--74, Mar./Apr. 2024. DOI: 10.1109/MSEC.2024.3352788.

\bibitem{backendal2024crypto}
M.~Backendal, H.~Davis, F.~G\"{u}nther, and K.~G.~Paterson, ``A formal treatment of end-to-end encrypted cloud storage,'' in \textit{Advances in Cryptology -- CRYPTO 2024}, Lecture Notes in Computer Science, Springer, Cham, 2024. DOI: 10.1007/978-3-031-68379-4\_2.

\bibitem{owasp2023}
OWASP Foundation, ``OWASP API Security Top 10 -- 2023,'' Open Web Application Security Project, 2023. [Online]. Available: https://owasp.org/API-Security/editions/2023/en/0x11-t10/.

\bibitem{owasp_session}
OWASP Foundation, ``Session Management Cheat Sheet,'' OWASP Cheat Sheet Series, 2023. [Online]. Available: https://cheatsheetseries.owasp.org/cheatsheets/Session\_Management\_Cheat\_Sheet.html.

\bibitem{cors2024}
MDN Web Docs Contributors, ``Cross-Origin Resource Sharing (CORS),'' Mozilla Developer Network, 2024. [Online]. Available: https://developer.mozilla.org/en-US/docs/Web/HTTP/Guides/CORS.

\bibitem{nist800-63b}
National Institute of Standards and Technology, ``NIST Special Publication 800-63B: Digital Identity Guidelines -- Authentication and Lifecycle Management,'' U.S. Dept. Commerce, Gaithersburg, MD, USA, 2020. [Online]. Available: https://pages.nist.gov/800-63-4/sp800-63b/session/.

\bibitem{apscheduler}
A.~Agronholm et al., ``APScheduler -- Advanced Python Scheduler,'' GitHub Repository and PyPI, 2024. [Online]. Available: https://github.com/agronholm/apscheduler.

\bibitem{celery}
Celery Project Contributors, ``Celery: Distributed Task Queue -- Official Documentation,'' v5.x, 2024. [Online]. Available: https://docs.celeryq.dev/.

\bibitem{google2013blog}
Google LLC, ``Plan your digital afterlife with Inactive Account Manager,'' \textit{Google Public Policy Blog}, Apr. 11, 2013. [Online]. Available: https://publicpolicy.googleblog.com/2013/04/plan-your-digital-afterlife-with.html.

\bibitem{google2024docs}
Google LLC, ``About Inactive Account Manager,'' \textit{Google Account Help Center}, 2024. [Online]. Available: https://support.google.com/accounts/answer/3036546.

\bibitem{flask2024}
Pallets Projects, ``Flask -- A Lightweight WSGI Web Application Framework,'' v3.1.x, Official Documentation, 2024. [Online]. Available: https://flask.palletsprojects.com/en/stable/.

\bibitem{react2024}
Meta Open Source / React Team, ``React -- The Library for Web and Native User Interfaces,'' Official Documentation, 2024. [Online]. Available: https://react.dev/.

\bibitem{sqlalchemy2024}
SQLAlchemy Authors, ``SQLAlchemy -- The Database Toolkit for Python,'' v2.0, Official Documentation, 2024. [Online]. Available: https://docs.sqlalchemy.org/en/20/.

\bibitem{doyle2023}
D.~T.~Doyle and J.~R.~Brubaker, ``Digital legacy: A systematic literature review,'' \textit{Proc. ACM Hum.-Comput. Interact.}, vol. 7, no. CSCW2, article 268, Oct. 2023. DOI: 10.1145/3610059.

\bibitem{futureme}
FutureMe.org, ``FutureMe -- Write a letter to your future self,'' [Online]. Available: https://www.futureme.org. [Accessed: Apr. 2026]

\bibitem{deadmansswitch}
Dead Man's Switch, ``Dead Man's Switch -- Automated email delivery service,'' [Online]. Available: https://www.deadmansswitch.net. [Accessed: Apr. 2026]

\bibitem{tiptap2024}
Tiptap GmbH, ``Tiptap -- The headless editor framework for web artisans,'' Official Documentation, v2.x, 2024. [Online]. Available: https://tiptap.dev/.

\bibitem{dompurify2024}
cure53, ``DOMPurify -- A DOM-only, super-fast, uber-tolerant XSS sanitizer for HTML, MathML and SVG,'' GitHub Repository, 2024. [Online]. Available: https://github.com/cure53/DOMPurify.

\bibitem{mediarecorder2024}
MDN Web Docs Contributors, ``MediaRecorder API,'' Mozilla Developer Network, 2024. [Online]. Available: https://developer.mozilla.org/en-US/docs/Web/API/MediaRecorder.

\end{thebibliography}

%=============================================================================
% APPENDICES
%=============================================================================
\newpage
\begin{appendices}

%-----------------------------------------------------------------------------
% APPENDIX A — SCREENSHOTS AND USER MANUAL OUTLINE
%-----------------------------------------------------------------------------
\section{Screenshots and User Manual}
\label{appendix:screenshots}

\subsection{Application Screenshots}

\begin{figure}[H]
  \centering
  \fbox{\parbox{0.85\textwidth}{\centering\vspace{3cm}[ADD SCREENSHOT — Home / Landing page]\vspace{3cm}}}
  \caption{Application Home Page}
  \label{fig:app_home}
\end{figure}

\begin{figure}[H]
  \centering
  \fbox{\parbox{0.85\textwidth}{\centering\vspace{3cm}[ADD SCREENSHOT — Registration form]\vspace{3cm}}}
  \caption{User Registration Screen}
  \label{fig:app_register}
\end{figure}

\begin{figure}[H]
  \centering
  \fbox{\parbox{0.85\textwidth}{\centering\vspace{3cm}[ADD SCREENSHOT — Capsule creation form (Time-based)]\vspace{3cm}}}
  \caption{Capsule Creation -- Time-Based Release}
  \label{fig:app_create_time}
\end{figure}

\begin{figure}[H]
  \centering
  \fbox{\parbox{0.85\textwidth}{\centering\vspace{3cm}[ADD SCREENSHOT — Capsule creation form (Event-based)]\vspace{3cm}}}
  \caption{Capsule Creation -- Event-Based Release}
  \label{fig:app_create_event}
\end{figure}

\begin{figure}[H]
  \centering
  \fbox{\parbox{0.85\textwidth}{\centering\vspace{3cm}[ADD SCREENSHOT — Delivery email received by recipient (HTML email format)]\vspace{3cm}}}
  \caption{Sample Capsule Delivery Email}
  \label{fig:delivery_email}
\end{figure}

\begin{figure}[H]
  \centering
  \fbox{\parbox{0.85\textwidth}{\centering\vspace{3cm}[ADD SCREENSHOT — Heartbeat ping email received by user]\vspace{3cm}}}
  \caption{Heartbeat Ping Email}
  \label{fig:heartbeat_email}
\end{figure}

\subsection{User Manual Summary}

\subsubsection{Getting Started}

\begin{enumerate}
  \item Navigate to the application URL in a modern web browser.
  \item Click \textbf{Register} and provide your full name, email address, and a password of at least 8 characters.
  \item After registration, you will be automatically logged in and redirected to your dashboard.
\end{enumerate}

\subsubsection{Creating Recipients}

\begin{enumerate}
  \item From the navigation bar, select \textbf{Recipients}.
  \item Click \textbf{Add Recipient} and enter the recipient's name, email address, and your relationship to them.
  \item Click \textbf{Save}. The recipient will appear in your recipients list.
\end{enumerate}

\subsubsection{Creating a Time Capsule}

\begin{enumerate}
  \item From the navigation bar, select \textbf{New Capsule} or click the \textbf{Create Capsule} button on the dashboard.
  \item Enter a title for your capsule.
  \item Compose your message using the rich text editor — use the toolbar to apply formatting (bold, italic, headings, lists, links). The editor accepts and stores formatted HTML.
  \item Optionally, click \textbf{Record Video} to open the in-browser video recorder. Grant camera and microphone permissions when prompted. Click \textbf{Start Recording}, record your message, then click \textbf{Stop}. Preview the recording and click \textbf{Attach This Recording} to attach it to the capsule.
  \item Select one or more recipients from your recipients list.
  \item Choose the release type: \textbf{Time-Based} (select a specific date and time) or \textbf{Event-Based} (triggered by inactivity detection or guardian verification).
  \item For Event-Based capsules, optionally enable \textbf{Requires Guardian Approval} and select one or more guardians.
  \item Click \textbf{Save as Draft} to save without scheduling, or \textbf{Schedule} to activate delivery.
  \item Optionally, upload additional file attachments from the Capsule Detail page after creation.
\end{enumerate}

\subsubsection{Managing Capsules}

\begin{itemize}
  \item \textbf{DRAFT} capsules can be edited and scheduled at any time.
  \item \textbf{SCHEDULED} capsules can be edited until delivery occurs.
  \item \textbf{SENT} and \textbf{CANCELLED} capsules are read-only.
\end{itemize}

\subsubsection{Guardian Verification}

For event-based capsules with guardian approval enabled:

\begin{enumerate}
  \item Navigate to the Capsule Detail page and scroll to the \textbf{Guardian Verification} section.
  \item Click \textbf{Request Verification} to dispatch verification emails to all assigned guardians. Each guardian receives a unique, tokenised link.
  \item Guardians navigate to their link (no account required) and see the capsule details. They select \textbf{Confirm Release} or \textbf{Deny Release}, optionally add a note, and submit.
  \item The capsule detail page shows a live audit log of all guardian responses. Click \textbf{Refresh} to update.
  \item When all guardians confirm, the capsule is delivered immediately to all recipients.
\end{enumerate}

\subsubsection{Heartbeat Confirmation}

\begin{itemize}
  \item Periodically, the system will send a \textit{``Are you still there?''} email to your registered address.
  \item Click the confirmation link in the email to reset the inactivity timer.
  \item If you do not respond within the expiry window, any event-based capsules you have configured will be automatically delivered to their designated recipients.
\end{itemize}

%-----------------------------------------------------------------------------
% APPENDIX B — API ENDPOINT TABLE
%-----------------------------------------------------------------------------
\section{Full API Endpoint Reference}
\label{appendix:api}

\begin{longtable}{|p{1cm}|p{4.5cm}|l|p{6cm}|}
\caption{Complete API Endpoint Reference}
\label{tab:api_full}\\
\hline
\textbf{Auth} & \textbf{Endpoint} & \textbf{Method} & \textbf{Description} \\
\hline
\endfirsthead
\hline
\textbf{Auth} & \textbf{Endpoint} & \textbf{Method} & \textbf{Description} \\
\hline
\endhead
\hline
\multicolumn{4}{r}{\textit{Continued on next page}}\\
\endfoot
\hline
\endlastfoot

No  & \texttt{/api/auth/register}          & POST   & Register a new user; returns JWT token \\
\hline
No  & \texttt{/api/auth/login}             & POST   & Authenticate user; returns JWT token \\
\hline
Yes & \texttt{/api/auth/me}                & GET    & Retrieve authenticated user profile \\
\hline
Yes & \texttt{/api/auth/refresh}           & POST   & Refresh a valid, non-expired token \\
\hline
Yes & \texttt{/api/recipients/}            & POST   & Create a new recipient \\
\hline
Yes & \texttt{/api/recipients/}            & GET    & List all recipients for authenticated user \\
\hline
Yes & \texttt{/api/recipients/<id>}        & GET    & Retrieve a specific recipient \\
\hline
Yes & \texttt{/api/recipients/<id>}        & PUT    & Update a recipient \\
\hline
Yes & \texttt{/api/recipients/<id>}        & DELETE & Delete a recipient \\
\hline
Yes & \texttt{/api/capsules/}              & POST   & Create a new capsule (encrypts message) \\
\hline
Yes & \texttt{/api/capsules/}              & GET    & List all capsules (message not returned) \\
\hline
Yes & \texttt{/api/capsules/<id>}          & GET    & Retrieve capsule with decrypted message \\
\hline
Yes & \texttt{/api/capsules/<id>}          & PUT    & Update capsule (re-encrypts if message changed) \\
\hline
Yes & \texttt{/api/capsules/<id>}          & DELETE & Delete capsule and all related records \\
\hline
Yes & \texttt{/api/capsules/<id>/attachments} & POST   & Upload a file attachment to a capsule \\
\hline
Yes & \texttt{/api/capsules/<id>/attachments} & GET    & List attachments for a capsule \\
\hline
Yes & \texttt{/api/capsules/<id>/attachments/<aid>} & DELETE & Delete a specific attachment \\
\hline
Yes & \texttt{/api/capsules/<id>/request-guardian-verification} & POST & Create GuardianVerificationRequest records and dispatch guardian emails \\
\hline
Yes & \texttt{/api/capsules/<id>/guardian-verifications} & GET & Return full guardian verification audit log for a capsule \\
\hline
No  & \texttt{/api/guardian/verify/<token>} & GET & Public: retrieve capsule info and current request status for a guardian token \\
\hline
No  & \texttt{/api/guardian/verify/<token>/respond} & POST & Public: guardian submits CONFIRM or DENY decision with optional notes \\
\hline
No  & \texttt{/api/health}                & GET    & Application health check (no auth required) \\
\hline

\end{longtable}

\noindent\textbf{Authentication:} All ``Yes'' endpoints require an \texttt{Authorization: Bearer <token>} header with a valid JWT. A 401 response is returned for missing, invalid, or expired tokens. A 403 response is returned for ownership violations.

\end{appendices}

%=============================================================================
% END OF DOCUMENT
%=============================================================================
\end{document}
